\documentclass[10pt,a4paper]{amsart}
\usepackage[margin=19mm]{geometry}
\usepackage{amsmath,amssymb,mathtools}
\usepackage[scr=rsfs,cal=euler]{mathalfa}
\usepackage{xfrac}
\usepackage[unicode,hidelinks,bookmarks=false]{hyperref}
\newcommand{\ie}{i.e.\ }
\newcommand{\cf}{cf.\ }
\newcommand{\resp}{respectively\ }
\newcommand{\Subsection}{\subsection}
\providecommand{\nobreakdash}{\nobreak\hbox{-}}
\setlength{\emergencystretch}{2em}
\multlinegap0pt
% ---------------------------------------------------------------------------
% Editorial AMS wrapper prelude. The theorem environments and the mathematical macros below are
% transcribed from the paper's own preamble (cached-originals/source0.tex lines 26--120); the AMS amsart
% class, the named format packages of the pinned TeX Live runtime and this editorial formatting are not
% part of the author's text. No mathematics is changed.
% ---------------------------------------------------------------------------
\usepackage{cleveref}
\newtheorem{theorem}{Theorem}[section]
\newtheorem{problem}{Problem}
\newtheorem{question}{Question}
\newtheorem{lemma}[theorem]{Lemma}
\newtheorem{corollary}[theorem]{Corollary}
\newtheorem{proposition}[theorem]{Proposition}
\newtheorem{claim}[theorem]{Claim}
\newtheorem{observation}[theorem]{Observation}
\newtheorem{conjecture}[theorem]{Conjecture}
\theoremstyle{definition}
\newtheorem{definition}{Definition}
\newtheorem{example}{Example}
\newtheorem{notation}{Notation}

\newcommand{\N}{\texttt{N}}
\newcommand{\E}{\mathbb{E}}


\theoremstyle{remark}
\newtheorem{remark}{Remark}
\sloppy
\newcommand{\red}[1]{{\color{red}#1}}
\newcommand{\blue}[1]{{\color{blue}#1}}
\newcommand{\pink}[1]{{\color{purple!85}#1}}

\newcommand{\Out}{\mathcal{O}}
\newcommand{\Sym}{\mathfrak{S}}
\newcommand{\NN}{\mathbb{N}}
\newcommand{\x}{\textbf{x}}
\newcommand{\North}{\texttt{N}}
\newcommand{\East}{\texttt{E}}
\newcommand{\DD}{\texttt{D}}
\newcommand{\UU}{\texttt{U}}
\newcommand{\tH}{\texttt{H}}
\newcommand{\calD}{\mathcal{D}}
\newcommand{\fd}{\Delta_{\ell} f}
\newcommand{\fdm}{\Delta_{m} f}
\newcommand{\SetL}[1]{S_{\Delta_{#1}}}
\newcommand{\mL}{\mathcal{L}}
\newcommand{\des}{\mathrm{des}}
\newcommand{\asc}{\mathrm{asc}}



\newcommand{\F}{\mathscr F}
\DeclareMathOperator{\PF}{PF}
\DeclareMathOperator{\PPF}{PPF}
\DeclareMathOperator{\disp}{dis} %changed form disp
\DeclareMathOperator{\Des}{Des}
\DeclareMathOperator{\tie}{tie}
\DeclareMathOperator{\Asc}{Asc}
\DeclareMathOperator{\Tie}{Tie}
\newcommand{\ER}{\mathbb{E}}
\newcommand{\PR}{\mathbb{P}}
\newcommand{\Luka}{\L ukasiewicz}
\newcommand{\LP}{\mathcal{L}}

\newcommand{\multinom}[2]{\left(\!\! \begin{array}{c} #1 \\ #2 \end{array} \!\!\right)}

\newcommand{\peh}[1]{\todo[size=\tiny,color=red!30]
{#1      \\ \hfill --- Pamela}}

\newcommand{\pamela}[1]{{\bf\color{purple}{[Pamela: #1}]}}
\newcommand{\selvi}[1]{{\bf\color{blue}{[Selvi: #1}]}}
\newcommand{\erin}[1]{{\bf\color{teal}{[Erin: #1}]}}
\newcommand{\mei}[1]{{\bf\color{orange}{[Mei: #1}]}}
\newcommand{\Kathryn}[1]{{\bf\color{cyan}{[Kathryn: #1}]}}
\newcommand{\tb}[1]{{\bf\color{blue}{#1}}}
\newcommand{\tr}[1]{{\bf\color{red}{#1}}}

\newcommand{\defterm}{\textbf}





\begin{document}
\title{On statistics of prime parking functions, Lukasiewicz paths, and quasisymmetric functions}
\author{Pamela E. Harris, Selvi Kara, Erin McNicholas, Kathryn Nyman and Mei Yin}
\date{Source: arXiv:2601.20770v2; distributed under CC BY 4.0}
\maketitle
\noindent\emph{Source, version and licence.} \emph{On statistics of prime parking functions, Lukasiewicz paths, and quasisymmetric functions}, by Pamela E. Harris, Selvi Kara, Erin McNicholas, Kathryn Nyman and Mei Yin. The source of this editable unit is the e-print of \textbf{version v2} of arXiv:2601.20770, https://arxiv.org/abs/2601.20770v2, whose material is distributed under the Creative Commons Attribution 4.0 International licence, http://creativecommons.org/licenses/by/4.0/, as recorded in the arXiv licence metadata for this paper and shown on that abstract page. The whole of the body below is version v2, the newest version of the paper. Full rights notice: licenses/arxiv-2601.20770-CC-BY-4.0.txt.\par\smallskip
\noindent\textit{Editorial scope.} Counted unit: original Theorem 6.4 of the article, the statement carried by the source environment \texttt{theorem} with source label \texttt{thm:Tie\_Set\_count\_mixed}, cached-originals/source0.tex lines 1276--1279, printed as \textbf{Theorem 6.4} exactly as in the article. It is delivered together with the authors' own complete proof as the single primary proof body, source0.tex lines 1287--1315 (29 lines): the \texttt{proof} environment that immediately follows the statement and counts the prime parking functions with the prescribed tie set. No proof marker is added and no mathematics is written, repaired or re-derived; the authors' notation and the printed slips of the source are kept literal. Necessary complete context is retained uncounted: source0.tex lines 1122--1127, Theorem 6.1, the corresponding single-parameter statement that the counted theorem generalizes and that the proof refers to for notation; the closure of the counted proof over references was computed and the further passages it needs are delivered with it. The article's own numbering is reproduced by explicit counter settings: the counter \texttt{theorem}, declared by \texttt{\textbackslash newtheorem\{theorem\}\{Theorem\}[section]} with lemma, corollary, proposition, claim, observation and conjecture sharing it, is set before each retained passage, so the counted statement prints as Theorem 6.4 and the cited statement as Theorem 6.1, exactly as in the article. Every \texttt{\textbackslash ref} and \texttt{\textbackslash Cref} inside every retained passage resolves inside this delivered file; because the article marks its cross-references with cleveref, the wrapper loads the same named package so that those commands typeset as in the source. No \texttt{\textbackslash cite} command occurs in any retained passage, so this unit delivers no bibliography entry; the article's own bibliography data file \texttt{bibliography.bib} is neither spliced into the bundled source nor redistributed. The retained passages contain no figure environment and no graphics-inclusion command, so no artwork is redistributed. The source's own class is AMS \texttt{amsart} (\texttt{\textbackslash documentclass[12pt,reqno]\{amsart\}}); the e-print ships no class or style file, so no publisher class is shipped, used or redistributed, and the wrapper is the same AMS \texttt{amsart} class with the article's own macros and theorem declarations transcribed from its preamble. Every declared body of this unit is an exact, unmodified span of source0.tex: no body is a bounded passage comparison, \texttt{omit} was not used anywhere, and consequently no fidelity receipt is asserted in delivery-evidence.json. The complete concatenated original source is bundled as sources/193547/source0.tex. Omitted from this editable unit and therefore absent from it are source0.tex lines 1--1105; 1110--1121; 1128--1275; 1280--1286; 1316--1338, that is the article's own preamble, title and author block and abstract, the introduction, the earlier sections on parking function shuffles, displacement enumerators and Lukasiewicz paths, the remaining quasisymmetric function results, and the article's bibliography data, all of which remain present in the bundled source but are not delivered here. The AMS \texttt{amsart} wrapper, the named format packages of the pinned TeX Live runtime and this editorial source, licence and scope paragraph are editorial formatting only.\par\smallskip
\setcounter{section}{6}
\setcounter{equation}{0}
\setcounter{definition}{6}
\setcounter{example}{7}
\setcounter{notation}{2}
\setcounter{problem}{0}
\setcounter{question}{0}
\setcounter{remark}{6}
\setcounter{theorem}{0}
\begin{theorem}\label{thm:Tie_Set_count}
    Given a set $S\subseteq[n-1]$, and fixed $0\leq \ell\leq n-2,$ the number of prime parking functions $\pi\in\PPF_n$ with $\SetL{\ell}(\pi)=S$ is  
    $$|\{\pi \in \PPF_n : \SetL{\ell}(\pi)=S\}| =(n-2)^{n-1-|S|}.$$
\end{theorem}

\setcounter{section}{6}
\setcounter{equation}{0}
\setcounter{definition}{6}
\setcounter{example}{7}
\setcounter{notation}{2}
\setcounter{problem}{0}
\setcounter{question}{0}
\setcounter{remark}{6}
\setcounter{theorem}{1}
\begin{theorem}\label{thm:quasisym}
For $n\geq 1$,
\begin{equation*}
\sum_{\pi\in\PPF_n}F_{n,\Tie(\pi)} =\sum_{i=1}^n (n-2)^{i-1}s_{(i,1^{n-i})}.
\end{equation*}
\end{theorem}

\setcounter{section}{6}
\setcounter{equation}{10}
\setcounter{definition}{6}
\setcounter{example}{8}
\setcounter{notation}{2}
\setcounter{problem}{0}
\setcounter{question}{0}
\setcounter{remark}{6}
\setcounter{theorem}{3}
\begin{theorem}\label{thm:Tie_Set_count_mixed}
    Let $m \neq \ell$. Given sets $S\subseteq[n-1]$ and $T\subseteq[n-1]$, and fixed $0\leq m, \ell\leq n-2,$ the number of prime parking functions $\pi\in\PPF_n$ with $\SetL{m}(\pi)=S$ and $\SetL{\ell}(\pi)=T$ is  
    $$|\{\pi \in \PPF_n : \SetL{m}(\pi)=S, \SetL{\ell}(\pi)=T\}| =(n-3)^{n-1-|S|-|T|}.$$
\end{theorem}

\setcounter{section}{6}
\setcounter{equation}{10}
\setcounter{definition}{6}
\setcounter{example}{8}
\setcounter{notation}{2}
\setcounter{problem}{0}
\setcounter{question}{0}
\setcounter{remark}{6}
\setcounter{theorem}{4}
\begin{proof}
We proceed as in the proof of \Cref{thm:quasisym}. Observe that
\begin{equation}
\label{eq:schur_mixed1}
\sum_{i=1}^n q^{n-i}(n-2)^{i-1} s_{(i,1^{n-i})}= \sum_{i=1}^n q^{n-i}(n-2)^{i-1} \sum_{S\subseteq [n-1], |S|=n-i}F_{n,S}.
\end{equation}
By change of variables $j=n-i$, we can rewrite \eqref{eq:schur_mixed1} as follows:
\begin{align*}
    \sum_{j=0}^{n-1} q^j (n-2)^{n-j-1} \sum_{S\subseteq [n-1], |S|=j} F_{n,S}
    =&\sum_{j=0}^{n-1}  \sum_{S\subseteq [n-1], |S|=j}  q^{|S|} F_{n,S} ~|\pi \in \PPF_n : \SetL{\ell}(\pi)=S|\\
    =&\sum_{\pi \in \PPF_n} q^{\fd(\pi)} F_{n,\SetL{\ell}(\pi)},
\end{align*}
where the first equality follows by \Cref{thm:Tie_Set_count}.

Similarly,
\begin{equation}
\label{eq:schur_mixed2}
\sum_{i=1}^n (q+n-3)^{i-1}s_{i,1^{n-i}}= \sum_{i=1}^n (q+n-3)^{i-1} \sum_{S\subseteq [n-1], |S|=n-i}F_{n,S}.
\end{equation}
By change of variables $j=n-i$, we can rewrite \eqref{eq:schur_mixed2} as follows:
\begin{align*}
    &\sum_{j=0}^{n-1} (q+n-3)^{n-j-1} \sum_{S\subseteq [n-1], |S|=j} F_{n,S}\\
    =&\sum_{j=0}^{n-1} \sum_{k=0}^{n-j-1} \binom{n-j-1}{k} q^k (n-3)^{n-j-k-1} \sum_{S\subseteq [n-1], |S|=j}  F_{n,S}\\
    =&\sum_{j=0}^{n-1} \sum_{k=0}^{n-j-1}   \sum_{\substack{S\subseteq [n-1] \\ |S|=j}} \sum_{\substack{T \subseteq [n-1]\setminus S \\ |T|=k}} q^{|T|} F_{n,S} ~|\pi \in \PPF_n : \SetL{m}(\pi)=S, \SetL{\ell}(\pi)=T| \\
    =&\sum_{\pi \in \PPF_n} q^{
\fd(\pi)} F_{n, \SetL{m}(\pi)},
\end{align*}
where the first equality uses the binomial expansion and the second equality follows by \Cref{thm:Tie_Set_count_mixed}.
\end{proof}

\end{document}
